% Volume VII: Formal Synthesis
% Part XIV. Dynamics, Geometry, and Holonomy
% Chapter 78: Nonintegrability and Institutional Path Dependence
% Spine: proof-spines/ch078-spine.md

\chapter{Nonintegrability and Institutional Path Dependence}
\label{ch:ch078}

Given the constrained distribution \(\mathcal D\subseteq T\mathcal M\), this chapter argues that institutional revision is often \textbf{nonintegrable}: the order of admissible operations matters. Hearing a public accusation before reviewing exculpatory material need not lead to the same endpoint as encountering the same materials in reverse order, even when the final dossier contains identical items. Sequence changes salience, credibility, procedural posture, and the meaning later evidence can acquire.

That claim gives formal expression to institutional path dependence. When \(\mathcal D\) is nonintegrable, there is no path-independent summary of inquiry; the route itself leaves residue in the final state. The chapter links this structure to \emph{narrative hysteresis} and to \textbf{distinction holonomy}: after traveling around a procedural loop, the institution may return to what looks like the same question while carrying a different classificatory orientation. With nonintegrability in place, the next chapter can represent that order-sensitivity geometrically through curvature.
